\documentclass[12pt]{article}
\usepackage[utf8]{inputenc}
\usepackage{amsmath, amssymb, amsthm}
\usepackage{geometry}
\usepackage{fancyhdr}

\geometry{letterpaper, margin=1in}
\setlength{\headheight}{15pt}

% Header setup
\pagestyle{fancy}
\fancyhf{}
\lhead{Math 311 Fall 2026}
\rhead{Jeffrey Lu}
\cfoot{\thepage}

\title{Homework Assignment 1 \\ \large Math 311, Fall 2026 \\ Rutgers University}
\author{}
\date{}

\begin{document}

\maketitle

\noindent \textbf{Exercise 1.2.5}
\begin{enumerate}
      \item [(a)]
            Suppose $x \in (A \cap B)^c$. Then, $ x \notin A \cap B $.
            So either $x \notin A$ or $x \notin B$.
            In other words, $x \in A^c$ or $x \in B^c$.
            Then, $ x \in A^c \cup B^c $.
            Therefore, $ (A \cap B)^c \subseteq A^c \cup B^c $.

      \item [(b)]
            Suppose $x \in A^c \cup B^c $. So $x \notin A$ or $x \notin B$.\
            Furthermore, $x$ couldn't be in both sets at the same time.
            So $x \notin A \cap B$, or $x \in (A \cap B)^c$.
            Therefore, $A^c \cup B^c \subseteq (A \cap B)^c$.
            Combining this with our result from part (a), we have that
            $A^c \cup B^c = (A \cap B)^c$.

      \item [(c)]
            Suppose $x \in (A \cup B)^c$. So $x \notin A \cup B$.
            Then, $x \notin A$ and $x \notin B$.
            In other words, $x \in A^c$ and $x \in B^c$.
            So $x \in A^c \cap B^c$.
            Thus, $(A \cup B)^c \subseteq A^c \cap B^c$.

            Next, suppose $y \in A^c \cap B^c$.
            Then, $y \notin A$ and $y \notin B$.
            That is, $y \notin A \cup B$.
            So $y \in (A \cup B)^c$.
            Then, $A^c \cap B^c \subseteq (A \cup B)^c $.

            Therefore, $A^c \cap B^c = (A \cup B)^c $.

\end{enumerate}

\noindent \textbf{Exercise 1.2.9}
\begin{enumerate}
      \item [(a)]
            $f^{-1}(A) = [-2, 2]$ and $f^{-1}(B) = [-1, 1]$.

            Notice that $A \cap B = [0, 1]$, and $f^{-1}(A \cap B) = [-1, 1]$.
            Similarly, $f^{-1}(A) \cap f^{-1}(B) = [-1,1]$ as well.
            Thus, $ f^{-1}(A \cap B) = f^{-1}(A) \cap f^{-1}(B) $

            Consider $A \cup B = [-1, 4]$. Then,
            \begin{align*}
                  f^{-1}(A \cup B) & = f^{-1}([-1, 4]) \\
                                   & = f^{-1}([0,4])   \\
                                   & = [-2, 2].
            \end{align*}
            Similarly, $f^{-1}(A) \cup f^{-1}(B) = [-2, 2]$.
            Thus, $ f^{-1}(A \cup B) = f^{-1}(A) \cup f^{-1}(B) $.
      \item [(b)]
            Let $x \in \mathbb{R}$. Then,
            \begin{align*}
                  x \in g^{-1}(A \cap B)
                   & \iff g(x) \in A \cap B                            \\
                   & \iff g(x) \in A \text{ and } g(x) \in B           \\
                   & \iff x \in g^{-1}(A) \text{ and } x \in g^{-1}(B) \\
                   & \iff x \in g^{-1}(A) \cap g^{-1}(B).
            \end{align*}
            Therefore,
            \[
                  g^{-1}(A \cap B)
                  = g^{-1}(A) \cap g^{-1}(B).
            \]

            Similarly,
            \begin{align*}
                  x \in g^{-1}(A \cup B)
                   & \iff g(x) \in A \cup B                           \\
                   & \iff g(x) \in A \text{ or } g(x) \in B           \\
                   & \iff x \in g^{-1}(A) \text{ or } x \in g^{-1}(B) \\
                   & \iff x \in g^{-1}(A) \cup g^{-1}(B).
            \end{align*}
            Therefore,
            \[
                  g^{-1}(A \cup B)
                  = g^{-1}(A) \cup g^{-1}(B).
            \]

\end{enumerate}

\noindent{} \textbf{Exercise 1.2.10}
\begin{enumerate}
      \item [(a)] False. Consider $a = b$.
            Then, for every $\epsilon > 0$, \[a < b + \epsilon.\]
            But, since $a = b$, the statement $a < b$ must be false.
            Thus, the if the and only if does not hold both ways.

      \item [(b)] False, same as the previous counter example.
      \item [(c)] True. Suppose
            $a \leq b$
            Since $\epsilon > 0$. It follows that
            \[
                  b < b + \epsilon.
            \]
            So,
            \[
                  a \leq b < b + \epsilon.
            \]
            That is,
            \[
                  a \leq b \implies a < b + \epsilon \mid \forall \epsilon > 0.
            \]

            Now, suppose that $a < b + \epsilon$ for every $\epsilon > 0$.
            For sake of contradiction, suppose $a > b$.
            Then $a - b > 0$. Choose
            \[
                  \epsilon = \frac{a-b}{2}.
            \]
            Then
            \begin{align*}
                  a < b + \dfrac{a-b}{2} = \dfrac{a+b}{2}
            \end{align*}
            But, since $a > b$,
            \[
                  a > \dfrac{a+b}{2}.
            \]
            Contradicting our previous inequality, so $a \leq b$.
            Therefore,
            \[
                  a < b + \epsilon \mid \forall \epsilon > 0\implies a \leq b.
            \].

            Hence,
            \[
                  a \leq b \iff a < b + \epsilon \mid \forall \epsilon > 0.
            \]

\end{enumerate}

\noindent{} \textbf{Exercise 1.2.13}
\begin{enumerate}
      \item [(a)] Let $S_n$ denote the statement that
            \[
                  (A_1 \cup A_2 \cup \dots \cup A_n)^c = A_1^c \cap A_2^c \cap \dots \cap A_n^c.
            \]

            (Base Case) Then $S_1$ says $(A_1)^c = A_1^c$, which is obviously true.

            (Inductive Step) We suppose $S_k$ for an arbitrary $k \in \mathbb{N}$. That is,
            \[
                  (A_1 \cup A_2 \cup \dots \cup A_k)^c = A_1^c \cap A_2^c \cap \dots \cap A_k^c.
            \]
            Consider
            \[
                  (A_1 \cup A_2 \cup \dots \cup A_k \cup A_{k+1})^c
            \]
            Then group the first $k$ sets,
            \[
                  ((A_1 \cup A_2 \cup \dots \cup A_k)\cup A_{k+1})^c.
            \]
            Apply 2-element De Morgan's Law from Exercise 1.2.5
            \[
                  (A_1 \cup A_2 \cup \dots \cup A_k)^c \cap A_{k+1}^c.
            \]
            Substitute $S_k$
            \[
                  A_1^c \cap A_2^c \cap \dots \cap A_k^c \cap A_{k+1}^c.
            \]
            That is,
            \[
                  (A_1 \cup A_2 \cup \dots \cup A_k \cup A_{k+1})^c = A_1^c \cap A_2^c \cap \dots \cap A_k^c \cap A_{k+1}^c.
            \]
            So $S_{k+1}$ holds. Therefore, by mathematical induction, $S_n$ is true for every $n \in \mathbb{N}$.
      \item [(b)]
            Let $B_i = (0, \frac{1}{i})$. Then for any $n \in \mathbb{N}$,
            \[
                  \bigcap_{i=1}^{n} B_i = (0, \frac{1}{n}).
            \]
            Since $n \in \mathbb{N}$, we can always
            find a larger natural number $n+1$.
            So $\frac{1}{n+1}$ will be inside the interval $(0, \frac{1}{n})$.
            Thus, the intersection is always nonempty.

            Now suppose
            \[
                  x \in \bigcap_{i=1}^{\infty} B_i.
            \]
            Then $ 0 < x < \frac{1}{i}$ for every $i \in \mathbb{N}$.
            Since $x > 0$, by the Archimedean Principle,
            there exists $i \in \mathbb{N}$ such that $\dfrac{1}{i} < x$.
            This contradicts $x < \dfrac{1}{i}$.
            Hence, no such $x$ exists.
            That is,
            \[
                  \bigcap_{i=1}^{\infty} B_i = \emptyset.
            \]
      \item [(c)]
            Suppose
            \[
                  x \in \left(\bigcup_{i=1}^{\infty} A_i\right)^c
            \]
            Equivalently,
            \[
                  x \notin \bigcup_{i=1}^{\infty} A_i
            \]
            That is, $x \notin A_i$ for every $ i \in \mathbb{N}$.
            Then $x \in A_i^c$ for every $ i \in \mathbb{N}$.
            Therefore,
            \[
                  x \in \bigcap_{i=1}^{\infty} A_i^c
            \]
            So
            \[
                  \left(\bigcup_{i=1}^{\infty} A_i\right)^c \subseteq \bigcap_{i=1}^{\infty} A_i^c
            \]

            Now, suppose
            \[
                  y \in \bigcap_{i=1}^{\infty} A_i^c
            \]
            Then $y \in A_i^c$ for every $ i \in \mathbb{N}$.
            Likewise, $y \notin A_i$ for every $ i \in \mathbb{N}$.
            That is,
            \[
                  y \notin \bigcup_{i=1}^{\infty} A_i
            \]
            So,
            \[
                  y \in \left(\bigcup_{i=1}^{\infty} A_i\right)^c
            \]
            Thus,
            \[
                  \bigcap_{i=1}^{\infty} A_i^c \subseteq \left(\bigcup_{i=1}^{\infty} A_i\right)^c
            \]
            Combining the two inclusions, we get
            \[
                  \left(\bigcup_{i=1}^{\infty} A_i\right)^c = \bigcap_{i=1}^{\infty} A_i^c.
            \]
\end{enumerate}

\noindent{} \textbf{Exercise 1.3.9}
\begin{enumerate}
      \item [(a)] Assume $\sup A < \sup B$.
            For the sake of contradiction,  assume $\sup A$ is an upper bound for $B$.
            Then,
            \[
                  \sup B \leq \sup A,
            \]
            contradicting $\sup A < \sup B$.
            Therefore, there exists $b \in B$ such that
            \[
                  b > \sup A.
            \]
            Furthermore,
            \[
                  a \leq \sup A < b \text{ for every } a \in A.
            \]
            Thus, $b$ is an upper bound for $A$.
      \item [(b)]
            Consider
            \[
                  A = B = (0, 1).
            \]
            Then
            \[
                  \sup A = \sup B = 1,
            \]
            and we have
            \[
                  \sup A \leq \sup B.
            \]

            But, no $b \in B$ can be an upper bound for $A$.
            For every $b \in B$, we can simply choose
            \[
                  a = \frac{b+1}{2}.
            \]
            Then $b < a < 1$, so $a \in A$ but $b < a$.
            Therefore, $b$ cannot be an upper bound for A.

\end{enumerate}

\noindent{} \textbf{Exercise 1.3.10}
\begin{enumerate}
      \item [(a)]
            Assume $A$ and $B$ are non-empty, disjoint sets with
            \[
                  A \cup B = \mathbb{R}
            \]
            and
            \[
                  a < b \text{ for every } a \in A, b \in B.
            \]
            Choose an arbitrary element $b_0 \in B$, then $b_0 > a$ for every $a \in A$.
            That is, $b_0$ is an upper bound for $A$.
            Since $A$ is non-empty and bounded above, the Axiom of Completeness guarantees us a least upper bound,
            call it
            \[
                  c = \sup A
            \]
            Since $c$ is an upper bound for $A$,
            \[
                  a \leq c \text{ for every } a \in A.
            \]
            Furthermore, because $c$ is the least upper bound on $A$ and every $b \in B$ is an upper bound for $A$,
            \[
                  c \leq b \text{ for every } b \in B.
            \]

            Therefore,
            \[
                  a \leq c \leq b
            \]
            for every $a \in A$ and $b \in B$.
      \item [(b)]
            Assume the Cut Property holds for $\mathbb{R}$.
            Let $E \subseteq \mathbb{R}$ be a nonempty set that is bounded above.
            Define
            \[
                  A = \{x \in \mathbb{R} : x \text{ is not an upper bound for } E\}
            \]
            and
            \[
                  B = \{x \in \mathbb{R} : x \text{ is an upper bound for } E\}.
            \]
            Then $A \cup B = \mathbb{R}$, $A \cap B = \emptyset$, and $B \neq \emptyset$ since $E$ is bounded above.
            Choose some $e \in E$. Then $e - 1 < e$. So $e - 1$ is not an upper bound on $E$.
            So $e - 1 \in A$ and $A \neq \emptyset$.
            Choose some $a \in A$ and $b \in B$.
            Since $a$ is not an upper bound on $E$, there exists some $x \in E$ such that
            \[
                  a < x.
            \]
            Since $b$ is an upper bound on $E$,
            \[
                  x \leq b.
            \]
            Thus,
            \[
                  a < x \leq b,
            \]
            so
            \[
                  a < b.
            \]
            Therefore the Cut Property holds. So there exists some $c \in \mathbb{R}$ such that
            \[
                  a \leq c \leq b
            \]
            for every $a \in A$ and $b \in B$.

            We want to prove that $c = \sup E$.

            Suppose, for sake of contradiction, that $c$ is not an upper bound on $E$.
            Then there exists some $x \in E$ such that
            \[
                  x > c.
            \]
            Choose
            \[
                  y = \frac{c+x}{2}.
            \]
            Then
            \[
                  c < y < x.
            \]
            Since $x \in E$ and $y < x$, $y$ cannot be upper bound on $E$, so $y \in A$.
            But the Cut Property says that every $a \in A$ satisfies $a \leq c$, so
            \[
                  y \leq c,
            \]
            contradicting $c < y$.
            Thus, $c$ must be an upper bound for $E$, so $c \in B$.
            By Cut Property,
            \[ c \leq b\] for every $b \in B$. That is, $c$ is the least upper bound of $E$, or
            \[
                  c = \sup E.
            \]
            Therefore, Cut Property implies the Axiom of Completeness.
      \item [(c)]
            Define
            \[
                  A = \{q \in \mathbb{Q} : q < \sqrt{2} \}
            \]
            and
            \[
                  B = \{q \in \mathbb{Q} : q > \sqrt{2} \}
            \]
            Then $A \cup B = \mathbb{Q}$, $A \cap B = \emptyset$, and $a < b$ for every $a \in A$ and $b \in B$.
            If the Cut Property were true in \(\mathbb Q\), there would be some \(c\in\mathbb Q\) such that
            \[
                  a \leq c \leq b
            \]
            for every $a \in A$ and $b \in B$.
            But the only number that can sit between all rationals below \(\sqrt2\) and all rationals above \(\sqrt2\) is
            \[
                  c=\sqrt2,
            \]
            and \(\sqrt2\notin\mathbb Q\).
            Therefore no such rational \(c\) exists, so the Cut Property fails in \(\mathbb Q\).
\end{enumerate}

\end{document}